% MY51990C
% Unterlassung der Reanimation
% 14.0
% 2013-11-30
:ref:`m-reanimation-unterlassung`
\Ca{Unter den folgenden Voraussetzungen darf (oder muss) eine
Reanimation unterlassen
werden:\index{Reanimation!Unterlassung
der}\index{Wiederbelebung|see{Reanimation}}}           
\begin{itemize}
  -   **Verwesung**
  -   **Fäulnis**
  -   **Mumifizierung**
  -   **Skelettierung**
  -   **Verletzungen, die nicht mit dem Leben
  vereinbar sind**
  -   **$2**\footnote{Unter
  \E{\enquote{Tierfraß}} versteht man den Verzehr eines Kadavers
  durch Kadaverfresser. Ein frischer Hundebiss oder
  Ähnliches fällt *nicht* darunter. Maden können
  sich im Gewebe von Lebenden einnisten und sind daher
  auch kein sicheres Todeszeichen.}
  -   (Verbindliche) **Patientenverfügung**\footnote{*Patientenverfügung*: Bei Notfallversorgungen darf **nicht nach einer Verfügung gesucht
  werden, wenn dadurch das Leben oder die Gesundheit des Patienten ernstlich
  gefährdet werden würde**.
  **Aber:** Wenn *zweifelsfrei* eine verbindliche
  Patientenverfügung vorliegt (z. B. weil eine betreuende
  Pflegekraft diese in der Zwischenzeit gefunden und die
  Verbindlichkeit bestätigt hat), muss dieser Folge
  geleistet werden.
Liegt eine *beachtliche Patientenverfügung* vor, so ist im
*Einzelfall* zu entscheiden: Das Wohl des Patienten bleibt
oberstes Gebot und die Verfügung soll in die
Entscheidung einfließen.}\index{Patientenverfügung!Unterlassung der
  Reanimation} (\FullPrettyRef{topic:patientenverfuegung}; 
  Bis zur zweifelsfreien Klärung der Situation muss
  jedenfalls eine Reanimation begonnen bzw. fortgeführt
  werden!)
  -   (Abbruch der Reanimationsmaßnahmen aufgrund
  von **Erschöpfung**)\footnote{Erschöpfung darf im
  professionellen Umfeld keine Rolle spielen, für eine
  entsprechende Ablösung ist bei Bedarf zu sorgen.}
  -   Unvereinbarkeit mit
  **Selbstschutz**\footnote{Es müssen jedoch alle Maßnahmen
  getroffen werden, die ohne Gefährdung des Selbstschutzes
  möglich sind. So wird zum Beispiel für Ersthelfer
  empfohlen, wenn
  eine Mund-zu-Mund-Beatmung nicht zumutbar ist, zumindest 
  eine Herzdruckmassage durchzuführen.}
\end{itemize}
